\documentclass{article}
\usepackage[utf8]{inputenc}
\usepackage{amsmath,amssymb}
\usepackage[most]{tcolorbox}
\usepackage{geometry}
\geometry{margin=2.5cm}

\newcommand{\step}[1]{\par\noindent\hangindent=3.5em\hangafter=1%
  \makebox[3.5em][l]{\textbf{#1.}}}
\newcommand{\steptag}[1]{\hfill{\small\textsf{[#1]}}}

\newtcolorbox{proofbox}[1]{
  colback=gray!5, colframe=gray!60,
  fonttitle=\bfseries, title={#1},
  breakable, enhanced,
  left=4pt, right=4pt, top=4pt, bottom=4pt,
  before skip=10pt, after skip=10pt
}

\begin{document}

\begin{proofbox}{There is a universal e\_idx\textasciicircum{}r > 0 such that, for every fin...}
\step{1} There is a universal e\_idx\textasciicircum{}r > 0 such that, for every finite-dimensional exact-unit epsilon\_r-C*-algebra with 0 <= epsilon\_r <= e\_idx\textasciicircum{}r and 1 < N = dim\_C calX < infinity, the scalar class breve-e = [J] is an isolated fixed point of the smooth breve-sigma, the vertical line iR*J is D-sigma\_J-invariant, ||D-breve-sigma\_\{breve-e\} + I|| < 1 in the quotient norm, and det(I - D-breve-sigma\_\{breve-e\}) > 0, so its local index is +1; more precisely, there is a quotient neighborhood calN of [J] such that if [U] in calN is fixed, choose a representative U\_0 close to J and c in U(1) with sigma(U\_0) = c*U\_0, choose a in U(1) with a\textasciicircum{}2 = c, and use sigma(a*U\_0) = conj(a)*sigma(U\_0) = a*U\_0: the two actual fixed lifts +-a*U\_0 lie in the J- and -J-isolation balls, hence equal J and -J, so [U] = [J]. \steptag{validated} \\[6pt]
\step{1.1} Fix W=(C\_rect,C\_ch,C\_pol,C\_grp,C\_path,C\_der,e\_rect,kappa\_ch,kappa\_pol,kappa\_der,delta\_*,epsilon\_*\textasciicircum{}r,e\_S1,r\_iso\textasciicircum{}W) from lem-stage1-polar-constant-ledger, (e\_iso\textasciicircum{}r,rho\_iso) from lem-stage1-uniform-inversion-isolation, e\_H\textasciicircum{}r from lem-stage1-quotient-left-inversion, and e\_quot\textasciicircum{}r from lem-stage1-quotient-manifold-package, and set e\_idx\textasciicircum{}r=min\{C\_rect*e\_S1,e\_iso\textasciicircum{}r,e\_H\textasciicircum{}r,e\_quot\textasciicircum{}r\}>0. If 0<=epsilon\_r<=e\_idx\textasciicircum{}r, put epsilon\_X=epsilon\_r/C\_rect; then epsilon\_X<=e\_S1, so clause (R) applies at delta=delta\_* and r=r\_iso\textasciicircum{}W, all hypotheses of (A\_2),(A\_4),(A\_5),(A\_7) hold, and C\_der*(epsilon\_r+r\_iso\textasciicircum{}W)<1. Clause (A\_2), lem-stage1-smooth-unitary-atlas, (A\_4), lem-stage1-smooth-polar-inverse, and lem-stage1-explicit-smooth-unitary-operations therefore synchronize a single smooth map sigma(U)=u\_\{delta\_*\}(U\textasciicircum{}dagger), with exactly the point values and first differential of the C\textasciicircum{}1 map in (A\_5)/(A\_7), satisfying sigma(J)=J and sigma(cU)=conj(c)*sigma(U); all three smallness-dependent quotient/isolation externals apply. \steptag{validated} \\[4pt]
\step{1.2} For 0<=epsilon\_r<=e\_idx\textasciicircum{}r and 1<N=dim\_C calX<infinity, lem-stage1-quotient-manifold-package and lem-stage1-quotient-left-inversion give M=breve-calU=calU\_e/U(1) as a compact orientable smooth manifold of real dimension N-1 and give a smooth descended self-map breve-sigma with breve-sigma([U])=[sigma(U)] and fixed point breve-e=[J]. For the quotient projection pi:calU\_e->M, the scalar orbit through J has tangent iR*J, so the differential of pi identifies T\_\{breve-e\}M with T\_J calU\_e/(iR*J), and the differential of any scalar-equivariant descended map is the operator induced on this quotient whenever its differential preserves iR*J. \steptag{validated} \\[6pt]
\step{1.2.1} Let e\_H\textasciicircum{}r>0 and e\_quot\textasciicircum{}r>0 be the universal cutoffs supplied respectively by lem-stage1-quotient-left-inversion and lem-stage1-quotient-manifold-package, and let the other positive universal quantities be those supplied by lem-stage1-polar-constant-ledger and lem-stage1-uniform-inversion-isolation. In the root existential construction choose, explicitly and independently of any sibling, e\_idx\textasciicircum{}r=min\{C\_rect*e\_S1,e\_iso\textasciicircum{}r,e\_H\textasciicircum{}r,e\_quot\textasciicircum{}r\}. Every entry of this finite minimum is positive, so e\_idx\textasciicircum{}r>0; moreover e\_idx\textasciicircum{}r<=e\_H\textasciicircum{}r and e\_idx\textasciicircum{}r<=e\_quot\textasciicircum{}r. Hence 0<=epsilon\_r<=e\_idx\textasciicircum{}r implies simultaneously epsilon\_r<=e\_H\textasciicircum{}r and epsilon\_r<=e\_quot\textasciicircum{}r. \steptag{validated} \\[4pt]
\step{1.2.2} Under 0<=epsilon\_r<=e\_idx\textasciicircum{}r and 1<N=dim\_C calX<infinity, the inequalities just proved meet the precise cutoff hypotheses of both externals: lem-stage1-quotient-manifold-package yields M=breve-calU=calU\_e/U(1) as a compact orientable smooth manifold of real dimension N-1, and lem-stage1-quotient-left-inversion yields the smooth descended map breve-sigma with breve-sigma([U])=[sigma(U)]. Since sigma(J)=J, breve-e=[J] is fixed. For the smooth quotient projection pi, the tangent to the U(1)-orbit t$|-\rangle$exp(it)J at J is iR*J; therefore Dpi\_J has kernel iR*J and identifies T\_\{breve-e\}M with T\_J calU\_e/(iR*J). Differentiating pi o sigma=breve-sigma o pi shows that whenever Dsigma\_J preserves iR*J, D-breve-sigma\_\{breve-e\} is exactly the endomorphism induced by Dsigma\_J on that quotient. \steptag{validated} \\[6pt]
\step{1.2.2.1} The missing quotient-differential facts follow directly from the already supplied smooth graph atlas and smooth scalar action, without assuming them from the quotient-manifold external.  The U(1)-action on calU\_e is free: if cU=U, then (c-1)U=0, and U has a right inverse, hence c=1.  Fix any U and use the graph chart chi\_U with omega\_U o Dchi\_U(0)=I\_\{iH\} and omega\_U(iU)=iJ.  Choose a real-linear functional ell:iH->R with ell(iJ)=1 and put K=ker ell.  For Theta\_U(t,k)=exp(it)chi\_U(k), defined near (0,0) in R x K, smoothness of the scalar action and chart gives omega\_U DTheta\_U(0,0)(s,k)=s*iJ+k, an isomorphism R x K -> iH; thus the finite-dimensional inverse-function theorem makes Theta\_U a diffeomorphism onto a neighborhood of U.  Because the action is free and U(1) is compact, after shrinking this neighborhood no phase outside the chosen small arc can return it to itself; therefore its saturation modulo U(1) is homeomorphic to an open subset of K, with quotient map represented in Theta\_U-coordinates by (t,k)$|-\rangle$k.  Repeating this construction at all U gives the canonical smooth quotient atlas: on an overlap, after multiplying one representative by the fixed phase matching the two centers, the transition is k$|-\rangle$pr\_K(Theta\_V\textasciicircum{}\{-1\}(Theta\_U(0,k))), hence is smooth.  Consequently pi:calU\_e->calU\_e/U(1) is a smooth submersion and ker Dpi\_U is the orbit tangent iR*U.  In particular ker Dpi\_J=iR*J and Dpi\_J induces T\_J calU\_e/(iR*J) congruent T\_[J](calU\_e/U(1)).  This canonical quotient has the same quotient topology as the manifold in lem-stage1-quotient-manifold-package, so the package supplies its compactness, dimension and topological orientability (equivalently smooth orientability); moreover the set map breve-sigma([U])=[sigma(U)] is smooth in this atlas because sigma is smooth and scalar-equivariant.  Finally pi o sigma=breve-sigma o pi; if Dsigma\_J preserves iR*J, the chain rule gives Dpi\_J Dsigma\_J=D-breve-sigma\_[J] Dpi\_J, and surjectivity of Dpi\_J says uniquely that D-breve-sigma\_[J] is the endomorphism induced by Dsigma\_J on T\_J calU\_e/(iR*J). \steptag{validated} \\[4pt]
\step{1.3} The covariance from lem-stage1-explicit-smooth-unitary-operations and sigma(J)=J imply, by differentiating sigma(exp(it)J)=exp(-it)J at t=0, that Dsigma\_J(iJ)=-iJ, hence iR*J is Dsigma\_J-invariant. At s=+1, clause (A\_7) of lem-stage1-polar-constant-ledger applies with delta=delta\_* and r=r\_iso\textasciicircum{}W. Here g\_J(0)=0 because f\_J(0)=0 and (A\_2) gives uniqueness, chi\_+(0)=J, and omega\_J is the inverse differential of chi\_+ at 0; since the smooth upgrades preserve point values and first derivatives, F\_+(A)=phi\_J\textasciicircum{}par(sigma(chi\_+(A))) is the graph-coordinate expression of sigma and DF\_+(0)=omega\_J*Dsigma\_J*omega\_J\textasciicircum{}\{-1\}. Thus ||DF\_+(0)+I\_\{iH\}||<=C\_der*(epsilon\_r+r\_iso\textasciicircum{}W)<1, and the invariant vertical line corresponds under omega\_J to iR*J because omega\_J(iJ)=iJ. \steptag{validated} \\[4pt]
\step{1.4} Let A=DF\_+(0) on the real normed space iH and L=iR*J. By the preceding invariant-line computation A(L) subseteq L, so A induces bar-A on iH/L; by the definition of the quotient norm, ||bar-A+I||<=||A+I||<1. Under the quotient tangent identification and graph-coordinate identification above, bar-A is D-breve-sigma\_\{breve-e\}; consequently ||D-breve-sigma\_\{breve-e\}+I||<1 in the quotient norm. \steptag{validated} \\[4pt]
\step{1.5} For any endomorphism T of a finite-dimensional real normed space with ||T+I||<1, set B=(T+I)/2. Then ||B||<1/2, every I-tB is invertible for 0<=t<=1 by the Neumann series, and continuity of the nonzero real determinant along t gives det(I-B)>0 from det(I)=1. Since I-T=2(I-B), det(I-T)=2\textasciicircum{}d det(I-B)>0, where d is the real dimension. Applying this to T=D-breve-sigma\_\{breve-e\} yields det(I-D-breve-sigma\_\{breve-e\})>0 (and in particular nonzero). \steptag{validated} \\[4pt]
\step{1.6} Shrink a quotient neighborhood calN of [J] so every class in calN has a representative U\_0 arbitrarily close to J. If [U\_0] is fixed by breve-sigma, then sigma(U\_0)=cU\_0 for some c in U(1). Continuity of sigma at J and ||J||=1 give |c-1|=||(c-1)J||<=||U\_0-J||+||sigma(U\_0)-J||, so after shrinking calN, c is arbitrarily close to 1. Label the two square roots of c as a and -a with a close to 1; shrinking once more makes aU\_0 lie in B\_\{rho\_iso\}(J) and -aU\_0 lie in B\_\{rho\_iso\}(-J). The covariance and a\textasciicircum{}2=c give sigma(aU\_0)=conj(a)cU\_0=aU\_0 and sigma(-aU\_0)=conj(-a)cU\_0=-aU\_0, so these are actual fixed lifts in the two prescribed isolation balls. \steptag{validated} \\[4pt]
\step{1.7} For calN from the phase-lift step, any fixed class [U] in calN has actual fixed representatives aU\_0 in B\_\{rho\_iso\}(J) and -aU\_0 in B\_\{rho\_iso\}(-J). Since epsilon\_r<=e\_idx\textasciicircum{}r<=e\_iso\textasciicircum{}r, lem-stage1-uniform-inversion-isolation forces aU\_0=J and -aU\_0=-J. Hence [U]=[U\_0]=[aU\_0]=[J]; thus breve-e=[J] is an isolated fixed point of breve-sigma, with precisely the representative/square-root conclusion stated in the root. \steptag{validated} \\[4pt]
\step{1.8} The smooth breve-sigma acts on the compact orientable manifold M, breve-e is isolated, and det(I-D-breve-sigma\_\{breve-e\})>0. Therefore lem-topology-local-index-sign applies and gives ind(breve-sigma,breve-e)=sgn det(I-D-breve-sigma\_\{breve-e\})=+1. \steptag{validated} \\[4pt]
\end{proofbox}

\end{document}
